\begin{table}[t]
\centering\footnotesize
\caption{Preregistered population-optimizer audit over 200 finite cyclic games. Regret differences are paired means with 95\% bootstrap intervals.}
\label{tab:synthetic-simplification}
\begin{tabular}{lcc}
\toprule
comparison & endpoint & result \\
\midrule
full vs. minimal & max policy $\ell_1$ & $0.0e+00$ \\
normalized vs. absorbed step & max policy $\ell_1$ & $2.9e-15$ \\
without Phase 0 & Nash-regret increase & $0.0187$ $[0.0165,0.0213]$ \\
noisy stage gate & Nash-regret change & $-0.0372$ $[-0.0407,-0.0340]$ \\
raw fixed step & strict-IR change & $-0.040$ $[-0.070,-0.015]$ \\
\bottomrule
\end{tabular}
\end{table}

\paragraph{Population optimizer audit.}
The full and minimal exact-gradient updates were identical, and absorbing the
sum of inverse-surplus weights into the step size changed the policy by at most
$2.9e-15$ in $\ell_1$. Removing Phase~0 increased
Nash-welfare regret by $0.0187$, while a stage-level acceptance test
reduced regret under frozen gradient noise by $0.0372$. Raw
inverse weights at the same nominal step reduced the strict-IR rate by
$4.0\%$. The scale-covariance and Pareto residuals were below
$2.4e-15$. These results support a minimal population update:
retain the feasibility warm start, normalize inverse-surplus weights or absorb
their sum into the step size, and omit post-warm-start clipping. Stage
acceptance remains an outer safeguard for inexact neural updates.

The simultaneous Hoeffding lower bound made no false certifications in this
study, compared with 1.09\% for the plug-in
rule at $n=128$, but its power was only 9.0\% at
$n=8192$. We therefore use confidence bounds as external certificates rather
than as an optimization ingredient.
